% ===== §12 The Beauty of Maintenance =====
\section{The Beauty of Maintenance: The Metaphysics of Repetition}
\label{p22:sec:maintenance}

\noindent
Perception discerns the field's direction and practice sends the path that way, but a field is not built once. It is built and rebuilt, turned and turned again, and the third face of the sensibility is the capacity to keep it accumulating across this turning rather than letting it dull into routine or spin in place. The beauty of maintenance is the temporal face, the holding of the good cycle in time, and it is the face on which the metaphysics of repetition bears, because to maintain a field is to repeat, and everything depends on what kind of repetition the maintaining is. This section develops maintenance and brings the metaphysics of repetition to the point the earlier parts prepared, the distinction between the repetition that gains and the repetition that turns in place, which is the distinction on which the difference between a sustained relation and a dead one rests.

\subsection{Maintenance, repetition, and the two repetitions}
\label{p22:sub:maintenance-two}

To maintain a field is to repeat: the same field is returned to, the same kinds of act performed again, the same shared life lived another day. This is why maintenance is the face on which repetition's metaphysics bears, and why the tradition's suspicion of the everyday, its association of the daily and the repetitive with the dull and the dead, must be met here. The suspicion has a basis, since one kind of repetition is exactly the dull and the dead, the mechanical re-enactment in which nothing is added and the same is merely done again, the wheel that turns without gaining. But this is one kind of repetition and not repetition as such. The other kind, the kind maintenance at its best performs, is the repetition that gains, in which the field is returned to and turned again with a difference, each turn carrying a finer perception and a more apt practice than the last, so that the repetition accumulates rather than recurs. The two are indistinguishable in any single return, since the same field is returned to in both, and they are told apart only across returns, by whether the field gains altitude or stays flat, which is the holonomy criterion brought to the temporal face of the sensibility.

\claim{Maintenance is repetition, and there are two repetitions. The bad repetition is the mechanical re-enactment in which the same is done again and nothing accumulates, the wheel whose holonomy is zero. The good repetition is the return that gains, in which the field is turned again with a difference, each turn carrying finer perception and apter practice, so that the field accumulates across returns. The two are indistinguishable within a single return and are told apart only across returns, by whether the field gains altitude. The beauty of maintenance is the capacity to make the return the gaining kind, to repeat so that the field accumulates rather than recurs, which is the holding of the good cycle in time.}

This distinction redeems the everyday from the tradition's suspicion without denying the suspicion its truth. The everyday is repetitive, and its repetition can be the dull mechanical kind that the suspicion fears; the dishcloth and the daily meal can be the wheel. But the repetitiousness of the everyday is not its condemnation, because the same daily return can be the gaining kind, the meal made again with a perception thickened by every previous making, the greeting repeated with a difference that accumulates. The everyday does not have to be escaped into the exceptional to become generative; it has to be repeated well, turned so that it gains, and the capacity to do this is the beauty of maintenance. This is why an aesthetic education of the everyday is possible at all, why the most repetitive and least spacious life is not excluded from generative cultivation: because the repetition that fills such a life can be the gaining kind, and the cultivation of maintenance is precisely the cultivation of the power to make it so.

\subsection{The craft tradition and the conditions of good repetition}
\label{p22:sub:maintenance-craft}

An old craft text states the conditions under which repeated making produces the good rather than the merely repeated, and it bears directly on the metaphysics of maintenance, because maintenance is a kind of making repeated. The text names four conditions that must meet for a made thing to be fine: the season has its times, the ground its airs, the material its beauty, and the craft its skill~\parencite{kaogongji}, and only when these four are joined can the work be good. The list is not a recipe but a statement of what good making depends on, and two of its features are exactly what the beauty of maintenance requires. The first is that good making is the meeting of a skill with conditions it does not control: the season and the ground and the material are given, not made, and the craft's skill consists in meeting them rightly, not in overriding them. This is the structure of maintenance as the paper has described it, the meeting of a trained perception with the state of a field it does not control, the skill lying in the apt response to given conditions rather than in their imposition. The maintainer of a field is the craftsman whose material is the field's state, whose skill is the perception of what that state affords, and whose good work is the meeting of skill with conditions rather than the forcing of conditions to a plan.

The second feature is that good making, on this account, is repeatable without being mechanical, because each making meets its own conditions of season and ground and material, which are never twice identical. The craftsman repeats the making, but the repetition gains rather than recurs, because each instance is the fresh meeting of the skill with conditions that have changed, the material this time not quite the material last time, the season turned. This is the good repetition stated as a craft ethic: the repeated making is not the mechanical reproduction of an identical object but the renewed meeting of a constant skill with varying conditions, which is exactly the gaining repetition of maintenance, the constant trained perception meeting the varying state of the field, each evening's walk a fresh meeting of the same skill with a season slightly turned. The craft tradition thus gives the beauty of maintenance an ancient model, the craftsman whose repeated making gains because it meets, each time, conditions it does not control, and whose skill is precisely the perception of what those conditions, this time, afford. To maintain a field well is to be such a craftsman, repeating the making of the field with a skill that meets its changing state, and the good repetition is the craftsman's repetition, gaining because it meets rather than imposes.

\subsection{Alienation and the ownership of the cycle's yield}
\label{p22:sub:maintenance-alienation}

The metaphysics of repetition must face the strongest case for the tradition's suspicion, which is not the dullness of the daily but its alienation, the condition in which repetition is genuinely deadening because it is not the repeater's own. The case is real and the section on the reproductive act met it from one side; it bears restating here as the limit of maintenance. There is a repetition that deadens not because it is repetition but because its yield is captured, the gesture repeated endlessly so that the holonomy it generates is siphoned to a power outside the relation, the worker's motion repeated for another's accumulation, and this repetition no cultivation of maintenance can redeem from within, because what deadens it is not the repeater's failure to turn the field well but the ownership of the field's yield by an order that takes the accumulation away. The beauty of maintenance presupposes that the holonomy the maintained field accumulates returns to the field, to the coupling that generates it, and where that presupposition fails, where the yield is captured, the failure is not aesthetic but structural, and its remedy is not a finer maintenance but a change in who owns the cycle's yield.

\claim{The repetition that genuinely deadens is the one whose holonomy is captured by an external power, not the one that is merely daily. A maintained field accumulates only if its yield returns to the coupling that generates it; where the yield is siphoned to an order outside the relation, the repetition deadens regardless of how well it is turned, and no cultivation of maintenance can redeem it, because what deadens it is structural and not aesthetic. The beauty of maintenance presupposes the return of the cycle's yield to the field, and the political economy of the field, taken up later, is the examination of when that presupposition holds and when it fails.}

This guards the metaphysics of repetition against a sentimentality it could otherwise fall into, the sentimentality of telling those whose repetition is captured that they need only turn their daily labour more beautifully to redeem it. They do not, and cannot, because the deadening of their repetition is not in their turning but in the ownership of its yield, and to tell them otherwise would be to ask them to find beauty in their own consumption. The beauty of maintenance is real and it is available in the everyday, but its availability presupposes a structural condition, that the field's yield returns to those who generate it, and where that condition fails the failure is to be named structurally and remedied structurally, not papered over with an exhortation to repeat more beautifully. The metaphysics of repetition thus carries its own political limit, the line past which the aesthetic question of how to repeat well gives way to the structural question of who owns what the repetition produces, and the honesty of the account requires marking that line rather than aestheticising across it.

\subsection{The dominant aspect within maintenance}
\label{p22:sub:maintenance-aspect}

Within the beauty of maintenance, as within the other two faces, a contradiction carries a dominant aspect that transforms with the field's state, and here the contradiction is between the freshening of the return and the steadiness of the return. To maintain a field well requires both that the return be fresh, carrying a difference that lets it gain, and that it be steady, reliable enough to be a return at all and not a series of unrelated novelties; a maintenance that was all freshness would not be the repetition of a field but its perpetual reinvention, and a maintenance that was all steadiness would be the mechanical wheel. In the ordinary case the dominant aspect is freshness, because the standing danger of maintenance is the slide into the mechanical, the dulling of the return into recurrence, and the work that matters is the keeping of the difference that lets the field gain. But the dominant aspect transforms: in a field thrown into disorder, where steadiness has been lost and the return itself is in doubt, the dominant aspect shifts to steadiness, the re-establishment of a reliable return without which no freshening has anything to freshen. Which aspect dominates is read from the field, the freshening in the field grown dull and the steadiness in the field grown chaotic, and the reading is once more the exercise of the sensibility whose temporal face maintenance is.

The interplay of freshness and steadiness repays one further observation, because it explains a feature of sustained relationships that is otherwise puzzling, the way the most enduring couplings often combine a deep reliability with a continual small surprise, and the way the loss of either is felt as a kind of death. A relation that has kept only its steadiness, that has become wholly reliable and wholly unsurprising, is felt to have died even though nothing is wrong, the death of the merely predictable, the wheel that turns so smoothly it is no longer felt to turn; this is the death by excess of steadiness, the dominant aspect having swung too far from freshness. A relation that has kept only its freshness, that is all surprise and no reliability, is felt to be exhausting and groundless, a series of novelties with no accumulated ground beneath them, unable to settle because nothing in it is steady enough to be returned to; this is the death by excess of freshness, the dominant aspect having swung too far the other way. The beauty of maintenance is the holding of the two in a balance that the field's state determines, steady enough to be a return and fresh enough to gain, and the art of it is the continual small adjustment of the balance as the field drifts toward one death or the other, adding freshness when the return dulls and adding steadiness when the return grows chaotic. This is why maintenance is the most temporally extended of the three faces and the one most easily mistaken for mere persistence: its work is invisible when it succeeds, felt only in its absence, the quiet continual rebalancing that keeps a long coupling from either of the deaths that threaten it, and that shows itself not as event but as the enduring aliveness of a relation that has been, though it does not announce it, continually maintained.
